% =============================================================
% preamble.tex — оформление отчётов по СТП БГУИР 01-2024
% =============================================================

% Шрифты: под pdflatex — Times через tempora + T2A,
% под xelatex/tectonic — системный Times New Roman (запасной вариант STIX Two Text)
\usepackage{iftex}
\ifPDFTeX
  \usepackage[utf8]{inputenc}
  \usepackage[T2A]{fontenc}
  \usepackage{tempora}
\else
  \usepackage{fontspec}
  \IfFontExistsTF{Times New Roman}
    {\setmainfont{Times New Roman}}
    {\setmainfont{STIX Two Text}}
  \IfFontExistsTF{DejaVu Sans Mono}
    {\setmonofont{DejaVu Sans Mono}[Scale=0.82]}
    {\setmonofont{DejaVuSansMono}[Extension=.ttf, BoldFont=*-Bold,
                                  ItalicFont=*-Oblique, Scale=0.82]}
\fi
\usepackage[russian]{babel}

% Поля: слева 30 мм, справа 15 мм, сверху и снизу 20 мм
\usepackage{geometry}
\geometry{left=30mm, right=15mm, top=20mm, bottom=20mm}

\usepackage{setspace}
\singlespacing

\pretolerance=1000
\tolerance=2000
\emergencystretch=3em

\usepackage{indentfirst}
\setlength{\parindent}{1.25cm}

% Номер страницы — снизу справа
\usepackage{fancyhdr}
\pagestyle{fancy}
\fancyhf{}
\fancyfoot[R]{\thepage}
\renewcommand{\headrulewidth}{0pt}
\fancypagestyle{plain}{%
  \fancyhf{}%
  \fancyfoot[R]{\thepage}%
  \renewcommand{\headrulewidth}{0pt}%
}

\usepackage{amsmath,amssymb}
\usepackage{graphicx}
\usepackage{float}
\usepackage{booktabs}
\usepackage{tabularx}
\usepackage{multirow}
\usepackage{longtable}
\usepackage{array}

\usepackage{caption}
\captionsetup[figure]{labelsep=endash, justification=centering,
                      font={normalsize}, name=Рисунок}
\captionsetup[table]{labelsep=endash, justification=raggedright,
                     singlelinecheck=false, font={normalsize}, name=Таблица}

\usepackage{titlesec}
\titleformat{\section}
  {\newpage\normalfont\fontsize{14}{18}\selectfont\bfseries}
  {\thesection}{0.5em}{\MakeUppercase}
\titlespacing*{\section}{1.25cm}{0pt}{14pt}
\titleformat{\subsection}
  {\normalfont\fontsize{14}{18}\selectfont\bfseries}
  {\thesubsection}{0.5em}{}
\titlespacing*{\subsection}{1.25cm}{14pt}{14pt}
\titleformat{\subsubsection}
  {\normalfont\fontsize{14}{18}\selectfont\bfseries}
  {\thesubsubsection}{0.5em}{}
\titlespacing*{\subsubsection}{1.25cm}{14pt}{14pt}

\newcommand{\sectioncenter}[1]{%
  \newpage
  \phantomsection
  \addcontentsline{toc}{section}{#1}
  \begin{center}\textbf{\MakeUppercase{#1}}\end{center}
  \vspace{14pt}
}

\usepackage{tocloft}
\renewcommand{\cftsecfont}{\normalfont}
\renewcommand{\cftsecpagefont}{\normalfont}
\renewcommand{\cftsubsecfont}{\normalfont}
\renewcommand{\cftsubsecpagefont}{\normalfont}
\renewcommand{\cftsecleader}{\cftdotfill{\cftdotsep}}
\renewcommand{\cftdotsep}{1}
% babel переопределяет заголовки после преамбулы, поэтому правим через \addto
\addto\captionsrussian{\renewcommand{\contentsname}{\hfill\textbf{СОДЕРЖАНИЕ}\hfill}}
\renewcommand{\cftaftertoctitle}{\hfill}

\usepackage{enumitem}
\setlist[itemize]{noitemsep, topsep=0pt, leftmargin=1.25cm, label=--}
\setlist[enumerate]{noitemsep, topsep=0pt, leftmargin=1.25cm}

% Листинги кода и вывода программы
\usepackage{listings}
\usepackage{fancyvrb}
\lstset{
  basicstyle=\ttfamily\footnotesize,
  frame=single,
  breaklines=true,
  showstringspaces=false,
  inputencoding=utf8,
  extendedchars=true,
  columns=fullflexible,
  keepspaces=true,
  captionpos=t
}
\ifPDFTeX
\lstset{
  literate={а}{{\cyra}}1 {б}{{\cyrb}}1 {в}{{\cyrv}}1 {г}{{\cyrg}}1
           {д}{{\cyrd}}1 {е}{{\cyre}}1 {ж}{{\cyrzh}}1 {з}{{\cyrz}}1
           {и}{{\cyri}}1 {й}{{\cyrishrt}}1 {к}{{\cyrk}}1 {л}{{\cyrl}}1
           {м}{{\cyrm}}1 {н}{{\cyrn}}1 {о}{{\cyro}}1 {п}{{\cyrp}}1
           {р}{{\cyrr}}1 {с}{{\cyrs}}1 {т}{{\cyrt}}1 {у}{{\cyru}}1
           {ф}{{\cyrf}}1 {х}{{\cyrh}}1 {ц}{{\cyrc}}1 {ч}{{\cyrch}}1
           {ш}{{\cyrsh}}1 {щ}{{\cyrshch}}1 {ъ}{{\cyrhrdsn}}1
           {ы}{{\cyrery}}1 {ь}{{\cyrsftsn}}1 {э}{{\cyrerev}}1
           {ю}{{\cyryu}}1 {я}{{\cyrya}}1 {ё}{{\cyryo}}1
           {А}{{\CYRA}}1 {Б}{{\CYRB}}1 {В}{{\CYRV}}1 {Г}{{\CYRG}}1
           {Д}{{\CYRD}}1 {Е}{{\CYRE}}1 {Ж}{{\CYRZH}}1 {З}{{\CYRZ}}1
           {И}{{\CYRI}}1 {Й}{{\CYRISHRT}}1 {К}{{\CYRK}}1 {Л}{{\CYRL}}1
           {М}{{\CYRM}}1 {Н}{{\CYRN}}1 {О}{{\CYRO}}1 {П}{{\CYRP}}1
           {Р}{{\CYRR}}1 {С}{{\CYRS}}1 {Т}{{\CYRT}}1 {У}{{\CYRU}}1
           {Ф}{{\CYRF}}1 {Х}{{\CYRH}}1 {Ц}{{\CYRC}}1 {Ч}{{\CYRCH}}1
           {Ш}{{\CYRSH}}1 {Щ}{{\CYRSHCH}}1 {Ъ}{{\CYRHRDSN}}1
           {Ы}{{\CYRERY}}1 {Ь}{{\CYRSFTSN}}1 {Э}{{\CYREREV}}1
           {Ю}{{\CYRYU}}1 {Я}{{\CYRYA}}1 {Ё}{{\CYRYO}}1
           {—}{{---}}1 {–}{{--}}1 {…}{{\ldots}}1
}
\fi

% Подписи листингов — в том же стиле, что таблицы: «Листинг X.Y — Название»
\captionsetup[lstlisting]{labelsep=endash, justification=raggedright,
                          singlelinecheck=false, font={normalsize}}
\renewcommand{\lstlistingname}{Листинг}
\AtBeginDocument{\numberwithin{lstlisting}{section}}

\lstdefinestyle{python}{language=Python, keywordstyle=\bfseries}
\lstdefinestyle{shell}{language=bash, keywordstyle=\mdseries}
\lstdefinestyle{plain}{language={}, basicstyle=\ttfamily\scriptsize}

% Абзац после листинга должен сохранять отступ 1,25 см (СТП 2.1.1)
\makeatletter
\AtBeginDocument{\let\@endpetrue\@endpefalse}
\makeatother

% Подпись к блоку verbatim: \captionof вне плавающего окружения
% глобально обнуляет \parindent, поэтому отступ восстанавливается вручную
\newcommand{\listingcaption}[2]{%
  \captionof{lstlisting}{#1}\label{#2}%
  \setlength{\parindent}{1.25cm}%
}

\usepackage{hyperref}
\hypersetup{colorlinks=true, linkcolor=black, urlcolor=blue, citecolor=black}

\numberwithin{figure}{section}
\numberwithin{table}{section}
\numberwithin{equation}{section}

% Титульный лист отчёта по лабораторной работе
\newcommand{\makelabtitle}[2]{%
  \begin{titlepage}
  \thispagestyle{empty}
  \begin{center}

  МИНИСТЕРСТВО ОБРАЗОВАНИЯ РЕСПУБЛИКИ БЕЛАРУСЬ

  \vspace{0.5em}
  Учреждение образования

  \vspace{0.5em}
  \textbf{<<Белорусский государственный университет информатики и радиоэлектроники>>}

  \vspace{2em}
  Факультет компьютерных систем и сетей

  \vspace{0.5em}
  Кафедра программного обеспечения информационных технологий
  \end{center}

  \vspace{3em}
  \begin{center}
  \textbf{ОТЧЕТ}

  \vspace{0.5em}
  по лабораторной работе №#1

  \vspace{0.5em}
  по дисциплине

  \vspace{0.5em}
  <<Информационные и технические средства защиты авторских прав>>

  \vspace{0.5em}
  <<#2>>
  \end{center}

  \vspace{5em}
  \begin{flushright}
  \begin{minipage}{0.5\textwidth}
  \textbf{Выполнил:}\\
  Лебедевич Артем Владимирович\\
  магистрант группы 556301

  \vspace{0.6em}
  \textbf{Преподаватель:}\\
  Ярмолик Вячеслав Николаевич\\
  доктор технических наук\\
  профессор кафедры ПОИТ
  \end{minipage}
  \end{flushright}

  \vfill
  \begin{center}Минск, 2026~г.\end{center}
  \end{titlepage}
}

% Титульный лист реферата
\newcommand{\makerefattitle}[1]{%
  \begin{titlepage}
  \thispagestyle{empty}
  \begin{center}

  МИНИСТЕРСТВО ОБРАЗОВАНИЯ РЕСПУБЛИКИ БЕЛАРУСЬ

  \vspace{0.5em}
  Учреждение образования

  \vspace{0.5em}
  \textbf{<<Белорусский государственный университет информатики и радиоэлектроники>>}

  \vspace{2em}
  Факультет компьютерных систем и сетей

  \vspace{0.5em}
  Кафедра программного обеспечения информационных технологий
  \end{center}

  \vspace{3em}
  \begin{center}
  \textbf{РЕФЕРАТ}

  \vspace{0.5em}
  по дисциплине

  \vspace{0.5em}
  <<Информационные и технические средства защиты авторских прав>>

  \vspace{1.5em}
  на тему

  \vspace{0.5em}
  \textbf{<<#1>>}
  \end{center}

  \vspace{5em}
  \begin{flushright}
  \begin{minipage}{0.5\textwidth}
  \textbf{Выполнил:}\\
  Лебедевич Артем Владимирович\\
  магистрант группы 556301

  \vspace{0.6em}
  \textbf{Преподаватель:}\\
  Ярмолик Вячеслав Николаевич\\
  доктор технических наук\\
  профессор кафедры ПОИТ
  \end{minipage}
  \end{flushright}

  \vfill
  \begin{center}Минск, 2026~г.\end{center}
  \end{titlepage}
}
